% Pursuing Stacks (cotes 134-2 to 134-8): archive cote 134-6, site page 61 (PDF page 62), typescript page 379.
% typescript. Draft compared: ps5.tex:4443-4519.
% AI-transcribed by Claude Opus 5.5 (claude-opus-5-5, adaptive thinking, effort high), Message Batches API, batch ps3 (msgbatch_01XdFmuVVBAcgc7EFxdtCqsj), 2026-11-12; not yet checked by eye.
\scanpage{134-6/61}
\mspage{379}

to equivalence)  each left hand factor to the category of functors
from the other functor to (Ab)  which commute with small direct limits
\xout{(}( much in the same way as the corresponding relationship between $A^{\wedge}$and
$B^{\wedge}$, with (Ab)  being replaced by (Sets))\xout{)} ,  and that accordingly the
functor (5) is equally fully faithful,  and extends to a fully faithful
functor from $B^{\wedge}_{ab}$ to  \uline{Homadd}$(A^{\wedge}_{ab},(Ab))$ ,  inducing in fact an equivalence
between $B^{\wedge}_{ab}$  and the full subcategory  \uline{Hom}$_{\hand{!}}(A^{\wedge}_{ab},(Ab))$ of \uline{Hom}$(A^{\wedge}_{ab},(Ab))$
made up by all functors  $A^{\wedge}_{ab}\hand{\to}(Ab)$  which commute to small direct limits\hand{.}
But for what we have in mind at present,  these niceties are not too re-%
levant yet it seems  -  all what matters is that an object L of Addinf(B)
defines an\xout{functorx} exact functor
\[ \overset{\hand{\sim}}{L} \ : \ A^{\wedge}_{ab} \longrightarrow (Ab) \ , \]
depending functorially on L,  in an infinitely additive way.  Thus,  as
noted in section 92  (but w\over{ith}{\unsure{here}} Addinf(B) \deleted{being}\inserted{was} replaced by the smaller
category Add(B),  which has turned out unsufficient for our purposes),
whenever we have a chain complex L.  in Addinf(B),  we get a correspondin\unsure{g}
functor
\[ \text{(7)} \qquad \overset{\hand{\sim}}{L}. \ : \ A^{\wedge}_{ab} \longrightarrow Ch.(Ab) \]
from $A^{\wedge}_{ab}$ to the category of chain complexes of (Ab),  which is moreover
an \uline{exact} functor.  Generalizing slightly the \xout{xxx} terminology introduced
in\inserted{section} 93,  \xout{xxxxx} where we restricted to chain complexes with components in
Add(B) rather than in Addinf(B),  we'll say that  $\overset{\hand{\sim}}{L}.$  is an \uline{abelianizator}
\uline{for}  A  ,  if the  following diagram commutes up to \xout{a}isomorphism :
\diagram{
\[ \text{(8)} \qquad
\begin{tikzcd}
A^{\wedge} \arrow[r] \arrow[d, "\begin{array}{c} Wh_A \\ \hand{(abelianization)} \end{array}"'] & Hot_A \arrow[r] & (Hot) \arrow[d, "\begin{array}{c} \hand{``absolute''} \\ \text{abelianization} \\ \text{functor} \end{array}"] \\
A^{\wedge}_{ab} \arrow[r, "\hand{\widetilde{L.}}"] & Ch.(Ab) \arrow[r] & (Hotab)
\end{tikzcd}
\qquad . \]
}
More accurately,  an abelianizator is a pair (L. ,\hand{$\lambda$} ) ,  where \hand{$\lambda$} is an
isomorphism of functors   $A^{\wedge} \longrightarrow (Hotab)$  making the diagram commute.
Here \xout{xxx} I like to view the abelianization functor
\[ \text{(9)} \qquad (Hot) \longrightarrow (Hotab) \overset{def}{=} D.(Ab) \]
as the one described directly in section (100) via integrators of arbi-%
trary modelizing objects i\xout{h}n (Cat),  without any reference to an auxiliary
test category such as $\Delta$ or the like. \hand{$\rfloor\!\lceil$}

The point\inserted{of (8)} is that  via an
"abelianizator" for A,  we want to be able to give  a)  a \uline{simultaneous}
handy expression,  in terms of ``computable'' chain complexes in (Ab),  of

\margin{61}
